% =====================================================================
% Dataset section for FinExam-10K.
%
% Structural choice worth stating up front: this section deliberately does
% NOT propose a hand-designed question-type taxonomy. Reviewer 2yXH on the
% companion paper rejected exactly that move as "unclear and non-mutually
% exclusive, lacks grounding in established reasoning literature". We instead
% report an empirical difficulty axis with measured reliability, which is a
% claim that can be checked rather than argued about, and keep the exams' own
% curriculum topics as the only categorical labels.
%
% \pending{} marks numbers awaiting the full answerability pass. The sampled
% estimates carry Wilson intervals and are reportable as they stand.
% =====================================================================
\newcommand{\pending}[1]{#1}   % 定稿前改成 \textcolor{red}{#1} 自查，然后删掉

\section{FinExam-10K}
\label{sec:dataset}

\textsc{FinExam-10K} is a benchmark of $10{,}198$ multiple-choice items drawn from the
professional examinations that license financial analysts and risk managers. Two properties
distinguish it from existing financial-reasoning corpora. First, difficulty is measured rather
than asserted: every item carries an empirical band derived from the frozen responses of
seventeen systems, and we report the reliability of that band. Second, half the corpus is
held out, so the benchmark can serve as a leaderboard rather than a one-time evaluation.

\subsection{Composition and Sources}

Items span all five stages of the two examinations, three CFA levels and two FRM parts, as
summarised in Table~\ref{tab:composition}. CFA items present three options and FRM items four,
so chance accuracy differs by track at $33.3$ and $25.0$ percent respectively, and all
cross-track comparisons in Section~\ref{sec:analysis} use chance-normalised accuracy.

Two collection streams feed the corpus. Mock and practice material published for examination
preparation forms the releasable half. Real examination items recalled and reconstructed by
candidates form the held-out half. We keep the two streams separate throughout, never mix
them within a release, and report every result on both halves so that readers can see whether
the two behave alike. They do: the two halves rank the seventeen systems at Spearman $0.990$,
while the public half is slightly the harder of the two, $16.8$ percent hard against $11.4$
percent, Mann-Whitney $p=3.2\times10^{-11}$. A benchmark whose public half is easier than its
hidden half would flatter every submission, so we note the direction explicitly.

\paragraph{From collection to release.}
Collection yielded $15{,}199$ items across the five stages. Two filters produce the released
corpus. We first drop every item that arrives without the provider's worked rationale, which
removes $1{,}484$ items, and then deduplicate, which removes a further $3{,}517$, leaving
$10{,}198$. The rationale filter is the reason coverage in the next paragraph is complete: a
worked rationale is a condition of entry rather than a property we were fortunate to find.

That condition is not neutral across stages and we report its effect rather than only its
benefit. Rationale availability is near total for CFA, $99.7$, $98.9$ and $100.0$ percent for
Levels I to III, but far lower for FRM, $68.4$ percent for Part~I and $61.6$ percent for
Part~II. The filter therefore removes about a third of the collected FRM material and almost no
CFA material, and the resulting $7{,}631$ to $2{,}567$ imbalance between the two examinations is
a consequence of that filter rather than of what was available to collect. Readers comparing
across examinations should treat the FRM subset as the survivors of a stricter selection.

\paragraph{Reference rationales.}
Every one of the $10{,}198$ items carries the examination provider's own worked rationale, with
no gaps in any stage or band. The rationales total $8.55$ million characters at a median of $752$
per item, and $29.6$ percent adjudicate each option in turn rather than only justifying the key,
in the form ``A Incorrect because \ldots\ B Correct because \ldots''. A quarter contain explicit
arithmetic. Most professional-examination corpora release the stem, the options and the keyed
letter alone, so this layer is what lets \textsc{FinExam-10K} support work that a bare
answer key cannot: verifying whether a system reached the right answer for the stated reason,
training and evaluating verifiers against a reference argument, and diagnosing which step of a
multi-step derivation a system lost. We use the rationales in this paper only to audit keyed
answers, and we do not place them in any model prompt.

Each item carries its examination, stage, curriculum topic, options, keyed answer, the provider
rationale, its empirical difficulty band and consensus pass rate, its release
split, and the integrity flags described in Section~\ref{sec:quality}. We do not impose a
question-type taxonomy of our own. The curriculum topics that the examination bodies already
publish are retained as the only categorical labels, because a taxonomy invented for a paper
invites the objection that its categories are neither exhaustive nor mutually exclusive, and
because the empirical axis below carries the information a hand-built taxonomy would be
proposing to capture.

\subsection{Empirical Difficulty}
\label{sec:difficulty}

The exams publish a level ordering, but that ordering is a curriculum sequence rather than a
statement about how hard an item is for a reasoning system. We therefore define
\texttt{difficulty\_v1} from model consensus. Seventeen frozen systems answer every item at
temperature $0$ under a fixed seed. The systems are partitioned into the ten proprietary ones
and the seven with open weights, the mean pass rate is taken
within each group, and the two group means are weighted equally so that band membership is not
decided by whichever group happens to be larger. Items with a consensus pass rate at or above
$2/3$ are easy, at or below $1/3$ are hard, and the remainder medium. This yields $6{,}578$ easy,
$2{,}183$ medium and $1{,}437$ hard items.

Three measurements support using this axis. Split-half reliability of the consensus score is
$0.805$ Spearman and $0.892$ after the Spearman-Brown correction. The worst leave-one-model-out
rank stability across the seventeen systems is $0.983$, so no single grader determines the
partition. And Cram\'er's $V$ against the exams' own level labels is only $0.174$, so the two
orderings are largely independent and the empirical axis is not a restatement of the published
levels. Appendix~\ref{app:difficulty-robustness} reports the same construction under a flat
seventeen-model average and under a three-group partition, and the appendix additionally rebuilds
the bands with DeepSeek-R1 and GPT-4o excluded from the panel, since those two systems appear as
backbones in the analysis of Section~\ref{sec:analysis}. Because the bands come from model
consensus, every band-conditioned result in the main text is scored leave-one-out, so no system
is graded against a partition it helped define.

\subsection{Quality Control and Known Defects}
\label{sec:quality}

\paragraph{Expert review.}
The corpus was reviewed by a four-person team with professional finance training, including
authors of this paper. Review covered the content of each item: that the stem, options and keyed
answer are faithful to the source, that the keyed answer agrees with the provider rationale, and
that the curriculum topic assignment is correct. Because a single pass yields no measurable
agreement, we additionally double-annotate a stratified sample of $400$ items and report Cohen's
$\kappa$ in Appendix~\ref{app:iaa}.

Content review and pipeline integrity are separate properties, and we report them separately. A
reviewer checking an item against its source sees the source, so a vignette or exhibit that was
lost during collection is visible to the reviewer and invisible in the shipped record. The
automated audits below therefore examine the released records themselves, independently of the
sources, and they surface defect classes that content review is not positioned to catch.

\paragraph{Structural integrity.}
Automated auditing finds no duplicate identifiers, no item whose keyed letter is absent from its
own option set, and no item missing a rationale. It surfaces four defect classes that we flag
rather than silently repair, and Table~\ref{tab:audit-outcomes} lists every count.

$111$ items contain two options that are identical after Unicode normalisation and whitespace
collapse, and in $80$ of those the keyed answer falls inside the duplicated pair, which makes the
item unscorable because a system selecting the other member gives textually the same answer and
is marked wrong. Five items contain an empty option and one presents only two options.

Deduplication during construction used exact matching, and a normalised comparison that also
folds case, punctuation and option order finds $1{,}192$ residual duplicate records in $1{,}148$
groups, $11.7$ percent of the corpus. Typical pairs differ only by a hyphen in
\emph{fixed-income}, by a trailing period, or by two options appearing in swapped order. This has
one consequence worth stating plainly: $104$ of those groups straddle the release boundary and
therefore expose $110$ held-out items, $2.2$ percent of the held-out half, to anyone reading
the public half. We ship the affected identifiers so that the leakage can be excluded, and future
releases will deduplicate under the normalised key. The corpus contains no contradictory key:
all $58$ groups whose keyed letter differs across members are explained by reordered options and
resolve to the same answer text.

\paragraph{Answerability.}
The larger defect is structural rather than typographic. Case-based examination items present a
shared vignette and an exhibit table, and both are frequently delivered as images that a
text-only collection pipeline does not capture. The result is items whose stem announces data it
never supplies, for example a stem reading ``Given the following ratings transition matrix,
calculate the two-period cumulative probability of default'' with no matrix, or one asking for a
named company's free cash flow with the company never introduced. Such items cannot be answered
from the text as shipped, and a system that fails them has demonstrated nothing about financial
reasoning.

A keyword detector for absent exhibits flags $732$ items, but it badly under-counts, because many
affected stems name no exhibit at all and simply reference an entity that was defined in the
missing vignette. We therefore label answerability directly for every item in the corpus. The
labeller was gated before use against $90$ items independently adjudicated by a stronger reader:
it agrees on $97.8$ percent, recovers $23$ of $23$ unanswerable items, and reaches a precision of
$92.0$ percent, so it misses no defect in the calibration set and over-flags mildly. Applied to
the full corpus it marks $2{,}573$ items unanswerable, $25.2$ percent, of which the keyword
detector had found only $714$.

\begin{center}
\small
\begin{tabular}{lrr}
\toprule
& \textbf{Not answerable} & \textbf{Rate} \\
\midrule
Easy   & $649/6{,}578$   & $9.9\%$ \\
Medium & $859/2{,}183$   & $39.3\%$ \\
Hard   & $1{,}065/1{,}437$ & $74.1\%$ \\
\midrule
CFA Level II & $1{,}209/2{,}215$ & $54.6\%$ \\
FRM Part I   & $214/1{,}652$     & $13.0\%$ \\
\midrule
Overall & $2{,}573/10{,}198$ & $25.2\%$ \\
\bottomrule
\end{tabular}
\end{center}

The dominant defect is an absent table, $1{,}550$ items, followed by an absent vignette, $846$,
an absent labelled object such as an undefined Statement~1, $142$, an option set belonging to a
different question, $27$, and a dangling reference, $8$. At the measured precision roughly two
hundred of the $2{,}573$ are expected to be false alarms, so the true count lies near $2{,}400$.

Because the defect concentrates in the hard band, every hard-band claim in
Section~\ref{sec:analysis} is reported twice, once over the full corpus and once over the
answerable subset, and the flag together with the detected missing referent ships with every
record so that users can reproduce or revise the partition.

We report this rather than quietly dropping the affected items for two reasons. The flag is
itself a usable research object, since the gap between answerable and unanswerable items measures
how a system behaves when its evidence is incomplete, which is a realistic failure mode. And a
benchmark that silently removes its inconvenient items cannot be audited by anyone else.

\subsection{Public Development and Sequestered Test Splits}

The $5{,}110$ mock and practice items are released in full. The $5{,}088$ recalled examination
items are held out and evaluated only through a submission endpoint. Splitting by collection
stream rather than at random means the public half cannot be used to infer the hidden half by
near-duplicate matching, which a random split over a corpus containing paraphrased variants would
not prevent.

Splitting by stream leaves the two halves unmatched in composition, and we state the differences
rather than average them away. The public half is $84.5$ percent CFA and $60.8$ percent Level I,
while the held-out half is $65.1$ percent CFA with Level I at $32.5$ percent and a
correspondingly larger FRM share. It also carries the higher defect rate, $33.3$ percent
unanswerable against $17.1$ percent.

That last difference reverses the apparent difficulty ordering, and the reversal is the practical
caveat for anyone tuning on this benchmark. Measured over the full corpus the public half looks
harder, with all seventeen systems scoring lower on it, $70.70$ against $75.28$. Measured over
the answerable subset the ordering flips: the public half is the easier of the two, $83.83$
against $81.63$, and no system scores lower on it. The apparent difficulty of the public half was
its defect rate rather than its content. Development scores computed on the answerable public
items should therefore be read as roughly two points optimistic relative to the held-out set. We
retain the stream-based split despite this, because it is what prevents a paraphrase in the
public half from revealing its counterpart in the held-out half, and a random split over a corpus
containing recalled variants of practice material would not. The seventeen systems still rank
almost identically on the two halves, Spearman $0.9902$, so the public half remains usable for
model selection even though its absolute level is offset.

\subsection{Versioning and Quarterly Refresh}
\label{sec:versioning}

\textsc{FinExam-10K} is maintained rather than frozen. We publish four versioned refreshes per
year, aligned with the February, May, August and November examination seasons. Each release adds
newly collected and expert-reviewed items where licensing permits, updates curriculum metadata,
repeats exact and semantic near-duplicate auditing, re-estimates \texttt{difficulty\_v1} against
the then-current model panel, and freezes a new public-development and held-out test manifest.

Not every stage is administered in every examination window, so the four checkpoints are
maintenance events rather than four complete replacements of all five subsets. Items are added to
the relevant stage when authorised material becomes available and passes the same review and
leakage-control procedures as the initial release.

Every leaderboard entry is bound to a dataset version, a model revision, a prompt configuration
and an evaluation date. New releases do not overwrite historical scores. The leaderboard retains
results for earlier versions and reports the current release separately, which lets the benchmark
track progress while limiting the staleness and repeated exposure that accumulate on a
permanently fixed test set. Because \texttt{difficulty\_v1} is defined against a specific model
panel, each release also records the panel used, and bands are never silently recomputed under an
existing version number.

\begin{table}[t]
\centering\small
\begin{tabular}{llrrr}
\toprule
\textbf{Exam} & \textbf{Stage} & \textbf{Items} & \textbf{Public} & \textbf{Held out} \\
\midrule
\multirow{3}{*}{CFA} & Level I   & 4{,}763 & 3{,}109 & 1{,}654 \\
                     & Level II  & 2{,}215 & 1{,}030 & 1{,}185 \\
                     & Level III & 653     & 179     & 474 \\
\midrule
\multirow{2}{*}{FRM} & Part I    & 1{,}652 & 516     & 1{,}136 \\
                     & Part II   & 915     & 276     & 639 \\
\midrule
\multicolumn{2}{l}{\textbf{Total}} & \textbf{10{,}198} & \textbf{5{,}110} & \textbf{5{,}088} \\
\midrule
\multicolumn{2}{l}{\textit{Difficulty band}} & & & \\
\multicolumn{2}{l}{\quad Easy}   & 6{,}578 & 3{,}164 & 3{,}414 \\
\multicolumn{2}{l}{\quad Medium} & 2{,}183 & 1{,}087 & 1{,}096 \\
\multicolumn{2}{l}{\quad Hard}   & 1{,}437 & 859     & 578 \\
\bottomrule
\end{tabular}
\caption{Composition of \textsc{FinExam-10K}. CFA items carry three options and FRM items four.
Bands are \texttt{difficulty\_v1}, defined in Section~\ref{sec:difficulty} from the consensus of
seventeen frozen systems rather than from the examinations' own level labels.}
\label{tab:composition}
\end{table}
